\newpage
\section[MapReduce: Simplified Data Processing on Large Clusters (37)]{MapReduce: Simplified Data Processing on Large Clusters (37) \cite{Dean2004MapReduce}}

\begin{enumerate}
    \item \textit{Category}: This paper is a description of a research prototype for the MapReduce programming model. The paper introduces concepts found in functional programming languages such as Haskell and Purescript, specifically map and reduce functions, and applies them to several operations that require computation. The problem that the MapReduce programming model aims to solve is the complexity related to the parallelization of multiple, simple computations for a task with a very large input data size to be done in a reasonable amount of time. It serves as an abstraction that seeks to hide away the messy details related to parallelization, fault-tolerance, data distribution, and load balancing in a library.
    \item \textit{Context}: Among the papers that are related to the selected paper are those about parallel and distributed computing systems, parallel programming and computation, and the Google cluster architecture. The theoretical bases used to analyze the problem in this paper are based on functional programming through the use of map and reduce functions, parallel computing, and data structures.
    \item \textit{Correctness}: The assumptions prove to be valid due to the nature of the problems presented in the paper as the computations involved that caused the problem all took in a mapping function, in which they are then reduced using a reducing function to get the derived data that they desire. Additionally, the added complexity regarding the parallelization of tasks due to a large input size also seems valid in order to not experience a slow execution time due to having sequential programs.
    \item \textit{Contributions}:
    The main contribution of this work is the design and large‐scale implementation of a simple yet powerful interface that enables:
    \begin{itemize}
        \item Automatic parallelization and distribution of large computations over commodity clusters,
        \item Transparent fault‐tolerance via re‐execution of failed tasks,
        \item Data‐locality optimizations to reduce network bottlenecks,
        \item And ease‐of‐use for programmers without extensive parallel or distributed computing expertise.
    \end{itemize}
    The paper also demonstrates real‐world impact at Google, reporting that hundreds of MapReduce programs are routinely run on thousands of machines. 
    \item \textit{Clarity}: The paper provided a clear introduction towards what it is mainly about. It also provided clear context on \textbf{why} the MapReduce model was made in its Introduction. For the conclusion of the paper, it managed to provide a satisfactory culmination of the meat of the paper by explaining why the model was proven to be a success and how it can be applied in the different fields of Google.
\end{enumerate}

    \textit{Refer to} \hyperref[fig2]{Appendix A2} for an \textbf{overview of the execution}


